\BourbakiAppendixExercises{习题}

\begin{enumerate}
\def\labelenumi{\arabic{enumi})}
\item
  给出一个体 D 上的可逆矩阵 A 的例子，它的转置不可逆；再给出 D 上的一个可逆矩阵 B 的例子，它的转置可逆但满足 \(\det^{t}B \neq \det B\)（考虑四元数体上的 2 阶方阵）。
\item
  设 D 是一个局部环；我们用 \(D^{*}\) 表示 D 的可逆元素所成的群，用 \(D_{ab}^{*}\) 表示 \(D^{*}\) 对其导群的商，用 \(\pi : D^{*} \to D_{ab}^{*}\) 表示典范同态。设 V 是一个维数 n 有限的自由右 D-模。我们用 \(\mathrm{B(V)}\) 表示 V 的基所成的集合，用 \(\Omega(\mathrm{V})\) 表示满足下列两个条件的映射 \(\omega : \mathrm{B(V)} \to \mathrm{D}_{\mathrm{ab}}^{*}\) 所成的集合：
\end{enumerate}

\begin{enumerate}
\def\labelenumi{(\roman{enumi})}
\item
  对 D* 中的 \(\lambda_1, \ldots, \lambda_n\) 与 \((v_1, \ldots, v_n) \in \mathrm{B}(\mathrm{V})\)，我们有 \(\omega(v_1 \lambda_1, \ldots, v_n \lambda_n) = \omega(v_1, \ldots, v_n)\pi(\lambda_1 \ldots \lambda_n)\)。
\item
  对 \(i \neq j\) 与 \(\lambda \in D\)，我们有 \(\omega(v_{1},\ldots,v_{i}+v_{j}\lambda,\ldots,v_{n}) = \omega(v_{1},\ldots,v_{n})\)。a) 设 e 是 V 的一个元素，W 是 V 的一个与 eD 互补的子模。设 \(\varphi \in \Omega(\mathrm{W})\)；证明存在唯一的元素 \(\omega\) 属于 \(\Omega(\mathrm{V})\)，满足
\end{enumerate}

\[
\omega (w _ {1}, \dots , w _ {n - 1}, e) = \varphi (w _ {1}, \dots , w _ {n - 1})
\]

对 W 的每个基 \((w_{1},\ldots ,w_{n - 1})\) 成立（把命题 1 的证明加以变通）。

\begin{enumerate}
\def\labelenumi{\alph{enumi})}
\setcounter{enumi}{1}
\item
  由此推出：给定两个元素 \(\omega\) 与 \(\omega'\)（属于 \(\Omega(\mathrm{V})\)），存在唯一的元素 t 属于 \(D_{ab}^{*}\)，使得我们有 \(\omega' = t\omega\)。
\item
  设 u 是 V 的一个自同构。证明存在 \(D_{ab}^{*}\) 的唯一元素，记作 \(\det u\)，使得我们有 \(\omega(u(v_{1}),\ldots,u(v_{n}))=(\det u)\omega(v_{1},\ldots,v_{n})\)（对 V 的每个基 \((v_{1},\ldots,v_{n})\) 与每个元素 \(\omega\)，后者属于 \(\Omega(\mathrm{V})\)）。映射 \(u\mapsto\det u\)（从 \(\mathbf{GL}(\mathrm{V})\) 到 \(D_{ab}^{*}\)）是一个群同态。
\end{enumerate}

\(d\) ) 取 \(\mathrm{V} = \mathrm{D}_d^n\)，于是 \(\mathbf{GL}(\mathrm{V})\) 与群 \(\mathbf{GL}_n(\mathrm{D})\)（\(n\) 阶可逆方阵所成的群）等同。证明映射 \(\operatorname{det}:\mathbf{GL}_n(\mathrm{D})\to \mathrm{D}_{\mathrm{ab}}^*\) 是唯一的在所有矩阵 \(B_{ij}(\lambda)\)（\(i\neq j\)，\(\lambda \in \mathrm{D}\)）上取值 1 且满足 \(\operatorname{det}(\operatorname{diag}(\lambda_1,\dots ,\lambda_n)) = \pi (\lambda_1\dots \lambda_n)\)（对 \(\lambda_1,\ldots ,\lambda_n\)，在 \(\mathrm{D}^*\) 中）的群同态。

\begin{enumerate}
\def\labelenumi{\alph{enumi})}
\setcounter{enumi}{4}
\tightlist
\item
  把第 4 目的例 2 至例 6 变通到这个框架中。
\end{enumerate}

¶ 3) 设 \(\mathrm{A} = \mathrm{A}_0 \oplus \mathrm{A}_1\) 是一个 \(\mathbf{Z}/2\mathbf{Z}\) 型的分次环。假设 \(\mathrm{A}_0\) 包含在 \(\mathrm{A}\) 的中心中，且 \(\mathrm{A}_1\) 的每个元素的平方为零。

\begin{enumerate}
\def\labelenumi{\alph{enumi})}
\item
  设 \(A_{1}^{2}\) 是 \(A_{0}\) 的由 \(A_{1}\) 中两个元素的乘积生成的 Z-子模；证明 \(A_{1}^{2}\) 是 \(A_{0}\) 的一个诣零理想，且 \(A_{1}^{2} \oplus A_{1}\) 是 A 的一个双边诣零理想。
\item
  设 \(p, q\) 是 \(\geqslant 0\) 的整数。用 \(A^{p|q}\) 表示自由分次右 A-模 \(A_d^p \oplus A_d(1)^q\)（II, §11, No.~2, 第 365 页，例 3），用 \(\mathbf{GL}(p|q, \mathrm{A})\) 表示 \(A^{p|q}\) 的（分次）（0 次）自同构所成的群。这样一个自同构在 \(A^{p|q}\) 的典范基下由一个矩阵 \(X = \begin{pmatrix} A & B \\ C & D \end{pmatrix}\) 表示，其中 \(A \in \mathbf{M}_p(\mathrm{A}_0)\)、\(B \in \mathbf{M}_{p,q}(\mathrm{A}_1)\)、\(C \in \mathbf{M}_{q,p}(\mathrm{A}_1)\)、\(D \in \mathbf{M}_q(\mathrm{A}_0)\)。证明矩阵 \(A\) 与 \(D\) 可逆（利用 \(a\) 与 VIII, 第 168 页的习题 16）。
\item
  对上述的 X，置 \(\operatorname{sdet} X = \det(A - BD^{-1}C) \det D^{-1}\)。证明 \(\operatorname{sdet} X\) 是 \(A_{0}\) 的一个可逆元素。
\end{enumerate}

\(d\) ) 证明每个矩阵 \(X\in \mathbf{GL}(p|q,\mathrm{A})\) 可以写成一个乘积

\[
X = \left( \begin{array}{c c} I & R \\ 0 & I \end{array} \right) \left( \begin{array}{c c} P & 0 \\ 0 & Q \end{array} \right) \left( \begin{array}{c c} I & 0 \\ S & I \end{array} \right).
\]

（取 \(R = BD^{-1}\)、\(P = A - BD^{-1}C\)、\(Q = D\)、\(S = D^{-1}C\)。）

\begin{enumerate}
\def\labelenumi{\alph{enumi})}
\setcounter{enumi}{4}
\tightlist
\item
  证明 sdet 是从 \(\mathbf{GL}(p|q,\mathrm{A})\) 到 \(A_{0}\) 的可逆元素所成的群的一个同态。（利用 d) 归结为证明等式 \(\operatorname{sdet}(XY)=\operatorname{sdet}X\operatorname{sdet}Y\)，当 X 形如 \(\binom{I}{S}\operatorname{sdet}(I)\)（其中 S 是一个初等矩阵）时。注意，对每个矩阵 \(R\in\mathbf{M}_{p,q}(\mathrm{A}_{1})\)，矩阵 RS（相应地 SR）的任意两个元素的乘积为零，并由此推出等式 \(\det(I-RS)=(\det(I-SR))^{-1}\)。
\item
  设 \(X = \begin{pmatrix} A & B \\ C & D \end{pmatrix}\) 是 \(\mathbf{GL}(p|q,\mathrm{A})\) 的一个元素；置 \(^{st}X = \begin{pmatrix} t_{A} & t_{C} \\ -t_{B} & t_{D} \end{pmatrix}\) 与 \(^{\pi}X = \begin{pmatrix} D & C \\ B & A \end{pmatrix}\)。设 \(Y \in \mathbf{GL}(p|q, A)\)；我们有 \(^{st}(XY) = ^{st}Y^{st}X\) 与 \(^{\pi}(XY) = ^{\pi}X^{\pi}Y\)。证明等式 \(\operatorname{sdet}X = \operatorname{sdet}^{st}X = (\operatorname{sdet}^{\pi}X)^{-1}\)（利用 d) 计算 \(\operatorname{sdet}^{\pi}X\)）。
\end{enumerate}

\begin{enumerate}
\def\labelenumi{\arabic{enumi})}
\setcounter{enumi}{3}
\tightlist
\item
  保留习题 3 的假设。设 \(\mathrm{M} = \mathrm{M}_0 \oplus \mathrm{M}_1\) 是一个有限生成的自由分次右 A-模；存在唯一的整数 \(p, q\)，使得 \(\mathrm{M}\) 同构于 \(\mathrm{A}^{p|q}\)。我们称 \(\mathrm{M}\) 的一个基是适配的，如果它的前 \(p\) 个向量是 0 次齐次的，后 \(q\) 个是 1 次齐次的。
\end{enumerate}

\begin{enumerate}
\def\labelenumi{\alph{enumi})}
\item
  设 u 是分次模 M 的一个自同构，\(M_{\mathcal{B}}(u)\) 是它关于 B 的一个适配基的矩阵。证明元素 sdet \(M_{\mathcal{B}}(u)\) 不依赖于 B 的选取。我们把它记作 sdet u。
\item
  用 \(u(1)\) 表示把自同态 u 视为分次模 M(1) 的自同构。若 B 是 M 的一个适配基，则交换 B 的前 p 个向量与后 q 个向量所得到的基 \(^{\pi}B\) 是 M(1) 的一个适配基，且我们有 \(M_{\pi\mathcal{B}}(u(1)) = ^{\pi}M_{\mathcal{B}}(u)\)（习题 3, f)）。由此推出我们有 \(\operatorname{sdet}(u(1)) = (\operatorname{sdet} u)^{-1}\)。
\item
  在每个分次左 A-模 N 上，我们通过置 \(ax = (-1)^{ij}xa\)（\(a \in A_i\)，\(x \in N_j\)）定义一个典范的分次右 A-模结构。特别地，我们把 M 的分次对偶 \(M^{*gr}\) 视为一个分次右 A-模，它同样同构于 \(A^{p|q}\)。设 u 是 M 的一个自同构。我们用 \({}^{st}u\) 表示 \(M^{*gr}\) 的由公式 \(\langle{}^{st}u(x^{*}), x\rangle = (-1)^{rs}\langle x^{*}, u(x)\rangle\)（\(x^{*} \in (\mathrm{M}^{*\mathrm{gr}})_r\)，\(x \in M_s\)）刻画的自同构。若 B 是 M 的一个适配基，则对偶基 \(B^*\) 是 \(M^{*gr}\) 的一个适配基，且我们有 \(M_{\mathcal{B}^*}({}^{st}u) = {}^{st}M_{\mathcal{B}}(u)\)（习题 3, f）。我们有 \(\operatorname{sdet}({}^{st}u) = \operatorname{sdet}(u)\)。
\end{enumerate}

